% part: intuitionistic-logic
% chapter: semantics
% section: introduction

\documentclass[../../../include/open-logic-section]{subfiles}

\begin{document}

\olfileid{int}{sem}{int}

\olsection{પરિચય}

અર્થવિજ્ઞાન વિના કોઈ તર્કપદ્ધતિનું સંતોષકારક વર્ણન થતું નથી;
સ્ફુરણાવાદી તર્ક પણ તેનો અપવાદ નથી. પ્રશિષ્ટ તર્ક માટે
!!{valuation}s પર આધારિત અર્થવિજ્ઞાન પ્રમાણભૂત છે, પરંતુ
સ્ફુરણાવાદી તર્ક માટે એકથી વધુ સ્પર્ધી અર્થવિજ્ઞાનો છે.
તેમાંનું કોઈ પણ સંયોજકોના અર્થનું સ્ફુરણાવાદી દૃષ્ટિએ
સંપૂર્ણપણે સ્વીકાર્ય વર્ણન આપે એ અર્થમાં પૂરતું સંતોષકારક નથી.

મોડલ તર્કના અર્થવિજ્ઞાન જેવું, !!{relational model}s પર
આધારિત અર્થવિજ્ઞાન કદાચ સૌથી વધુ પ્રચલિત છે. તેમાં
!!{propositional variable}s જગતોને સોંપવામાં આવે છે અને
આ જગતોને સુલભતા સંબંધ જોડે છે. આ સંબંધ હંમેશા આંશિક
ક્રમ હોય છે, એટલે સ્વવાચક, વિસંમિત અને પરંપરિત હોય છે.

સહજ રીતે આ જગતોને જ્ઞાનની અવસ્થાઓ અથવા ``પુરાવાની
પરિસ્થિતિઓ'' માની શકાય. અવસ્થા~$w'$ એ~$w$માંથી ત્યારે અને
માત્ર ત્યારે જ સુલભ છે જ્યારે આપણને જાણીતું છે તે મુજબ
$w'$ એ જ્ઞાનની શક્ય (ભાવિ) અવસ્થા હોય, એટલે~$w$માં
જે જાણીતું છે તેની સાથે સુસંગત હોય. એક વાર કોઈ વિધાનની
ખબર પડે, તો તે અજાણ્યું બની શકતું નથી: $!A$ એ~$w$માં
જાણીતું હોય અને~$Rww'$ હોય, તો $!A$ એ~$w'$માં પણ
જાણીતું હોય. તેથી સુલભતા સંબંધની દૃષ્ટિએ ``જ્ઞાન''
એકદિશવર્ધી છે.

જ્ઞાનસંબંધી તર્કની જેમ ``$!A$ જાણીતું છે''નો અર્થ
``બધા જ્ઞાનસંબંધી વિકલ્પોમાં સત્ય છે'' એવો લઈએ, તો
બધા વિકલ્પોમાં $!A$ અને~$!B$ બંને જાણીતાં હોય ત્યારે
$!A \land !B$ એ~$w$માં જાણીતું છે. પરંતુ જ્ઞાન
એકદિશવર્ધી છે અને $R$ સ્વવાચક છે; તેથી $!A \land !B$
એ~$w$માં ત્યારે અને માત્ર ત્યારે જ જાણીતું હોય છે જ્યારે
$!A$ અને $!B$ બંને $w$માં જાણીતાં હોય. એ જ કારણસર
$!A \lor !B$ એ~$w$માં ત્યારે અને માત્ર ત્યારે જ
જાણીતું હોય છે જ્યારે બેમાંથી ઓછામાં ઓછું એક જાણીતું હોય.
તેથી $\land$ અને $\lor$ માટે સંયોજકોની સત્ય-શરતો
પ્રશિષ્ટ તર્ક જેવી જ છે.

પરંતુ શરતી સંયોજકની સત્ય-શરતો પ્રશિષ્ટ તર્કથી જુદી છે.
$!A \lif !B$ એ~$w$માં ત્યારે અને માત્ર ત્યારે જ જાણીતું
હોય છે જ્યારે $Rww'$ ધરાવતી કોઈ પણ અવસ્થા~$w'$માં $!A$
જાણીતું હોય અને સાથે $!B$ જાણીતું ન હોય એવું ન બને.
આ શરત~$w$માં $!A$ અજાણ્યું હોય અથવા $!B$ જાણીતું
હોય તે શરત જેવી નથી. કારણ કે~$w$માં $!A$ કે $!B$
બેમાંથી કોઈ જાણીતું ન હોય, છતાં $Rww'$ ધરાવતી ભાવિ
જ્ઞાનસંબંધી અવસ્થા~$w'$ હોઈ શકે; એ $w'$માં $!A$
જાણીતું થાય અને $!B$ જાણીતું ન થાય એવું બની શકે છે.

$\lnot !A$ની ખબર આપણને ત્યારે જ હોય જ્યારે કોઈ શક્ય
ભાવિ જ્ઞાનસંબંધી અવસ્થામાં~$!A$ની ખબર પડી શકે તેમ ન હોય.
વિચાર એ છે કે $!A$ જાણવું શક્ય હોય, તો કોઈ શક્ય ભાવિ
અવસ્થામાં~$!A$ જાણીતું બને. $\lfalse$ની ખબર પડી શકતી
નથી, એટલે એવી ભાવિ અવસ્થામાં~$!A$ની ખબર હોય પણ
~$\lfalse$ની ન હોય.

આ અર્થઘટનમાં વર્જિત મધ્યનો નિયમ નિષ્ફળ જાય છે. કેટલાંક
$!A$ અત્યારે આપણને જાણીતાં નથી, પરંતુ આગળ જતાં જાણી
શકીએ. એવા !!a{formula}~$!A$ માટે $!A$ અને
~$\lnot !A$ બંને અજાણ્યાં છે, તેથી $!A \lor \lnot !A$
પણ જાણીતું નથી. પરંતુ, દાખલા તરીકે, $\lnot(!A \land
\lnot !A)$ જાણીતું છે: $!A$ અને~$\lnot !A$ બંનેની
ખબર હોય એવી કોઈ ભાવિ અવસ્થા શક્ય નથી. આપણને હાલ $!A$
કે $\lnot !A$ની ખબર છે કે નહીં તેનાથી સ્વતંત્ર રીતે આ
વાત જાણીએ છીએ.

!!^{relational model}s જ સ્ફુરણાવાદી તર્ક માટે ઉપલબ્ધ
એકમાત્ર અર્થવિજ્ઞાન નથી. સ્થલરચનાત્મક અર્થવિજ્ઞાન બીજું
છે: તેમાં વિધાનોને સ્થલરચનાત્મક અવકાશમાં ખુલ્લા ગણો
તરીકે અર્થઘટિત કરીએ છીએ, અને સંયોજકોને આ ગણો પરની
સંક્રિયાઓ તરીકે (દાખલા તરીકે, $\land$ને છેદને અનુરૂપ)
અર્થઘટિત કરીએ છીએ.

\end{document}
